\section{System Model and Problem Formulation}

\subsection{System Model}

\subsubsection{Network Scenario}

本研究考慮一個 OneWeb-like 的 LEO 衛星通訊系統，並聚焦於 NGSO--GSO coexistence 情境下的下行干擾管理問題。系統中每顆 LEO 衛星配置 $16$ 個 spot beams，並於每一個 time slot 中在衛星覆蓋範圍內隨機生成地面 users。設 $\mathcal{S}(t)$ 為當前 time slot $t$ 中可見之 LEO 衛星集合，$\mathcal{B}_s$ 為衛星 $s$ 的 beam 集合，且 $|\mathcal{B}_s|=16$。

在每一個 time slot 中，每位 user 僅能由一顆衛星的一個 beam 服務。設 $\mathcal{U}(t)$ 為當前時槽中的 user 集合，並以二元變數 $x_{u,s,b}(t)$ 表示 user--beam association。若 user $u$ 由衛星 $s$ 的 beam $b$ 服務，則 $x_{u,s,b}(t)=1$，否則為 $0$。由於每位 user 僅能對應到單一 serving satellite 與單一 beam，因此 association 必須滿足唯一性限制。

此外，系統中設有一個受保護之 GSO ground station (GS)。由於 LEO 衛星在特定時間可能與 GS 呈現近 inline 幾何關係，因此當某顆 LEO 衛星接近 GS 正上方時，其下行發射較容易造成較高的 aggregate EPFD。本文將該時刻對 GS 造成最大干擾風險的 LEO 衛星定義為 critical satellite，記為 $s^{\star}(t)$。同時，考慮 critical satellite 同軌道上的北側與南側鄰近衛星，分別記為 $s^{\mathrm{N}}(t)$ 與 $s^{\mathrm{S}}(t)$，作為後續 relay / reassociation 的候選 helper satellites。

\subsubsection{User Association and Service Model}

設由衛星 $s$ 的 beam $b$ 所服務之 user 集合為
\begin{equation}
\mathcal{U}_{s,b}(t)=\{u\in\mathcal{U}(t)\mid x_{u,s,b}(t)=1\}.
\end{equation}

則 beam 上的 user 數定義為
\begin{equation}
N_{s,b}(t)=\sum_{u\in\mathcal{U}(t)} x_{u,s,b}(t).
\end{equation}

每位 user 的需求速率記為 $D_u(t)$。若所有 users 具有相同需求，則可簡化為固定常數 $D$，例如本研究設定為 $25$ Mbps。

由於每位 user 僅能由單一衛星之單一 beam 服務，因此需滿足
\begin{equation}
\sum_{s\in\mathcal{S}(t)}\sum_{b\in\mathcal{B}_s} x_{u,s,b}(t)\le 1,
\qquad \forall u\in\mathcal{U}(t).
\end{equation}

此外，若 user $u$ 不在 beam $b$ 的可服務覆蓋區 $\mathcal{C}_{s,b}(t)$ 內，則不可建立 association，即
\begin{equation}
x_{u,s,b}(t)=0,
\qquad \forall u\notin\mathcal{C}_{s,b}(t).
\end{equation}

\subsubsection{Load-Aware Power Allocation}

每顆衛星 $s$ 具有固定最大總發射功率 $P_s^{\max}$。此外，考量實際多波束衛星酬載架構中，各 beam 的發射能力通常仍受限於 beam-level RF chain、power amplifier 或 EIRP 上限，因此本文進一步對每個 beam 設定最大可分配功率，記為 $P_{s,b}^{\max}$。

在初始 user association 完成後，本文依據 active beams 上的 user 數比例分配衛星總功率，以反映較高負載的 beam 應獲得較多傳輸功率。設 $\mathcal{B}_s^{\mathrm{act}}(t)$ 為衛星 $s$ 在當前 time slot 的 active beam 集合，則 beam $b$ 的初始分配功率可表示為
\begin{equation}
\begin{aligned}
P_{s,b}^{(0)}(t)
&=
\frac{N_{s,b}(t)}
{\sum_{b'\in\mathcal{B}_s^{\mathrm{act}}(t)} N_{s,b'}(t)}
\,P_s^{\max},
\end{aligned}
\end{equation}
其中 $b\in\mathcal{B}_s^{\mathrm{act}}(t)$。

為避免將過高功率不切實際地集中於單一 beam，本文進一步考慮每個 beam 的最大功率限制，因此 beam 的最終分配功率表示為
\begin{equation}
P_{s,b}(t)=\min\!\left(P_{s,b}^{(0)}(t),\,P_{s,b}^{\max}\right),
\qquad b\in\mathcal{B}_s^{\mathrm{act}}(t).
\end{equation}

上述配置提供一個 load-aware 的初始功率分配結果。然而，最終 beam power allocation 除需滿足衛星總功率限制與 beam-level 最大功率限制外，仍須進一步滿足對受保護 GSO ground station 的 aggregate EPFD constraint，因此後續仍需執行 power control 或其他資源調整機制。

\subsubsection{SINR and User Satisfaction Model}

在本研究中，user 的接收品質由 serving beam 所提供之期望訊號、同頻干擾以及接收雜訊共同決定。與 full frequency reuse 不同的是，本文採用 OneWeb-like 8-channel frequency assignment。因此，user 所受到的 co-channel interference 並非來自所有 active beams，而僅來自與其 serving beam 使用相同 channel 的 beams。

設 $\mathcal{K}$ 為可用 channel 集合，且 $|\mathcal{K}|=8$。若 user $u$ 由衛星 $s$ 的 beam $b$ 服務，則該 beam 所使用之 channel 記為 $k_{s,b}\in\mathcal{K}$。則 user $u$ 的接收 SINR 可定義為
\begin{equation}
\mathrm{SINR}_u(t)=
\frac{S_u(t)}
{I_u^{\mathrm{intra}}(t)+I_u^{\mathrm{inter}}(t)+N_u(t)}.
\end{equation}

其中 $S_u(t)$ 為期望訊號功率，$I_u^{\mathrm{intra}}(t)$ 與 $I_u^{\mathrm{inter}}(t)$ 分別表示 intra-satellite 同頻干擾與 inter-satellite 同頻干擾，$N_u(t)$ 為接收端雜訊功率。

\paragraph{(1) Desired signal power}
對於由衛星 $s$ 的 beam $b$ 所服務之 user $u$，其期望訊號功率可表示為
\begin{equation}
S_u(t)=
x_{u,s,b}(t)\,
\frac{P_{s,b}(t)\,G^{\mathrm{tx}}_{s,b\to u}(t)\,G_u^{\mathrm{rx}}}
{L_{s,u}(t)}.
\end{equation}

其中 $P_{s,b}(t)$ 為 beam $b$ 的發射功率，$G^{\mathrm{tx}}_{s,b\to u}(t)$ 為 beam 朝 user 方向之發射增益，$G_u^{\mathrm{rx}}$ 為 user 端接收增益，$L_{s,u}(t)$ 為衛星 $s$ 與 user $u$ 之總傳播損失。

本文將總傳播損失表示為
\begin{equation}
L_{s,u}(t)=L_{s,u}^{\mathrm{fs}}(t)\,L_{s,u}^{\mathrm{gas}}(t),
\end{equation}
其中 $L_{s,u}^{\mathrm{fs}}(t)$、$L_{s,u}^{\mathrm{gas}}(t)$ 分別代表 free-space path loss 與 atmospheric gas loss。對於 free-space path loss，其線性形式可寫為
\begin{equation}
L_{s,u}^{\mathrm{fs}}(t)=\left(\frac{4\pi d_{s,u}(t)}{\lambda}\right)^2,
\end{equation}
其中 $d_{s,u}(t)$ 為衛星 $s$ 與 user $u$ 之間的距離，$\lambda$ 為載波波長。

\paragraph{(2) Intra-satellite interference}
由於同一顆衛星上的不同 beams 可能重複使用相同 channel，因此 user 可能受到同一顆 serving satellite 上其他同頻 beams 的干擾。本文將此類干擾定義為 intra-satellite interference，其表達式為
\begin{equation}
\begin{aligned}
I_u^{\mathrm{intra}}(t)
&=
\sum_{\substack{b'\in\mathcal{B}_s\\ b'\neq b}}
\frac{P_{s,b'}(t)\,G^{\mathrm{tx}}_{s,b'\to u}(t)\,G_u^{\mathrm{rx}}}
{L_{s,u}(t)}.
\end{aligned}
\end{equation}

\paragraph{(3) Inter-satellite interference}

除同星同頻干擾外，user 亦可能受到其他可視衛星之同頻 beams 所造成之干擾。本文將此類干擾定義為 inter-satellite interference，其表達式如下。

\begin{equation}
I_u^{\mathrm{inter}}(t)=
\sum_{\substack{s'\in\mathcal{S}(t)\\ s'\neq s}}
\;
\sum_{b'\in\mathcal{B}_{s'}}
\frac{
P_{s',b'}(t)\,G^{\mathrm{tx}}_{s',b'\to u}(t)\,G_u^{\mathrm{rx}}
}{
L_{s',u}(t)
}.
\end{equation}

上式表示，來自其他可視衛星且與 serving beam 使用相同 channel 的 beams，皆可能對 user 造成 co-channel interference。

\paragraph{(4) SINR expression}

綜合上述定義，user 的接收 SINR 可表示如下。

\begin{equation}
\mathrm{SINR}_u(t)=
\frac{S_u(t)}
{I_u^{\mathrm{intra}}(t)+I_u^{\mathrm{inter}}(t)+N_u(t)}.
\end{equation}

上式顯示，在 OneWeb-like 8-channel 架構下，user 的接收品質不僅取決於 serving beam 的功率與幾何位置，亦同時受到 intra-satellite 與 inter-satellite beams 干擾影響。因此，beam power allocation 與 channel-aware user association 將直接影響系統之整體服務表現。

\paragraph{(5) Instantaneous capacity and actual service rate}
在求得 user 的 SINR 後，可依 Shannon formula 計算其 instantaneous capacity：
\begin{equation}
C_u(t)=B_u\log_2\!\bigl(1+\mathrm{SINR}_u(t)\bigr).
\end{equation}

其中 $B_u$ 為 user 所使用之有效頻寬。

由於同一 beam 內可能同時服務多位 users，本文以 TDMA 描述 beam 內資源共享。設 user $u$ 由衛星 $s$ 的 beam $b$ 服務，且該 beam 上共有 $N_{s,b}(t)$ 位 users，則其實際服務速率可表示為
\begin{equation}
R_u(t)=\frac{C_u(t)}{N_{s,b}(t)}.
\end{equation}

\paragraph{(6) User satisfaction}

為評估 user 實際服務品質是否滿足其需求，本文定義 satisfaction 如下。

\begin{equation}
\mathrm{sat}_u(t)=\min\!\left(\frac{R_u(t)}{D_u(t)},\,1\right).
\end{equation}

其中，$D_u(t)$ 為 user $u$ 之需求速率。當 $\mathrm{sat}_u(t)=1$ 時，表示 user 需求已被完全滿足；當 $\mathrm{sat}_u(t)<1$ 時，表示實際服務速率不足以完全支撐其需求。此 satisfaction metric 將作為後續 problem formulation 與 performance evaluation 的主要指標之一。

\subsubsection{Beam-Level EPFD Model for 16-Beam Satellite}

在 NGSO--GSO coexistence 情境下，除需考慮 user link quality 外，亦必須確保 LEO 衛星對受保護 GSO ground station (GS) 所造成之 aggregate EPFD 不超過法規限制。由於本研究中每顆 LEO 衛星配置 16 個 beams，且不同 beams 對 GS 的方向、離軸角與發射功率可能不同，因此單顆衛星對 GS 的干擾不能以單一總功率近似，而應以 beam-level contribution 個別計算後，再於線性域加總。

\paragraph{(1) Beam-level PFD toward the protected GS}
考慮衛星 $s$ 的 beam $b$，其朝受保護 GS 方向所形成之 power flux density (PFD) 可表示為
\begin{equation}
\mathrm{PFD}_{s,b}^{\mathrm{GS}}(t)=
\frac{P_{s,b}(t)\,G^{\mathrm{tx}}_{s,b\to g}(t)}
{4\pi d_{s,g}^{2}(t)}.
\end{equation}

其中 $G^{\mathrm{tx}}_{s,b\to g}(t)$ 為 beam 朝 GS 方向之發射增益，$d_{s,g}(t)$ 為衛星 $s$ 與受保護 GS 之間的距離。

\paragraph{(2) Beam-level EPFD contribution}

EPFD 除考慮到衛星發射至 GS 所形成之 PFD 外，亦需進一步考慮 GS 接收天線對 LEO 方向的離軸接收增益，並相對於其最大接收增益進行正規化。故 beam 對 GS 的 EPFD contribution 可表示如下。

\begin{equation}
\mathrm{EPFD}_{s,b}(t)=
\frac{P_{s,b}(t)\,G^{\mathrm{tx}}_{s,b\to g}(t)}
{4\pi d_{s,g}^{2}(t)}
\cdot
\frac{G_g^{\mathrm{rx}}(\theta_{g\to s}(t))}
{G_{g,\max}^{\mathrm{rx}}}
\cdot
\frac{1}{B_{\mathrm{ref}}}.
\end{equation}

其中，$G_g^{\mathrm{rx}}(\theta_{g\to s}(t))$ 為 GS 天線朝衛星 $s$ 方向之接收增益，$G_{g,\max}^{\mathrm{rx}}$ 為 GS 接收天線最大增益，$B_{\mathrm{ref}}$ 為 EPFD 評估所採用之 reference bandwidth。

\paragraph{(3) Angular dependence of beam gain}
對於衛星 $s$ 的 beam $b$，GS 相對於該 beam boresight 的離軸角記為 $\phi_{s,b\to g}(t)$。因此，beam 在 GS 方向之發射增益可視為離軸角的函數，即
\begin{equation}
G^{\mathrm{tx}}_{s,b\to g}(t)=G^{\mathrm{tx}}_{s,b}\!\bigl(\phi_{s,b\to g}(t)\bigr).
\end{equation}

由於同一顆衛星上的 16 個 beams 各自對應不同覆蓋方向，因此即使它們共享同一顆衛星平台，對 GS 的離軸角仍不相同，進而使各 beam 的 EPFD contribution 顯著不同。這也是本文採用 beam-level EPFD calculation，而非以單一衛星總功率近似的主要原因。

\paragraph{(4) Per-satellite EPFD aggregation over 16 beams}

由於每顆衛星具有 16 個 beams，因此衛星 $s$ 對 GS 的總 EPFD contribution 應為所有 active beams 貢獻之線性總和。設 $\mathcal{B}_s^{\mathrm{act}}(t)$ 為衛星 $s$ 在當前 time slot 的 active beam 集合，則衛星 $s$ 的總 EPFD contribution 可表示如下。

\begin{equation}
\mathrm{EPFD}_{s}^{\mathrm{sat}}(t)=
\sum_{b\in\mathcal{B}_s^{\mathrm{act}}(t)} \mathrm{EPFD}_{s,b}(t).
\end{equation}

由於本研究中每顆衛星固定配置 16 個 beams，因此亦可寫為

\begin{equation}
\mathrm{EPFD}_{s}^{\mathrm{sat}}(t)=
\sum_{b\in\mathcal{B}_s} \mathrm{EPFD}_{s,b}(t),
\end{equation}

其中對於未啟動或未分配功率之 beam，可視為 $P_{s,b}(t)=0$。

\paragraph{(5) Aggregate EPFD from multiple visible satellites}
在當前 time slot 中，所有可視 LEO 衛星對受保護 GS 所形成之 aggregate EPFD 可表示為
\begin{equation}
\mathrm{EPFD}_{\mathrm{agg}}(t)=
\sum_{s\in\mathcal{S}(t)} \mathrm{EPFD}_{s}^{\mathrm{sat}}(t).
\end{equation}

等價地，也可先在單顆衛星內部對 16 個 beams 加總，再對所有可視衛星做加總：
\begin{equation}
\mathrm{EPFD}_{\mathrm{agg}}(t)=
\sum_{s\in\mathcal{S}(t)}\sum_{b\in\mathcal{B}_s} \mathrm{EPFD}_{s,b}(t).
\end{equation}

由於 EPFD 需於線性域累加，因此在進行多 beam、多衛星干擾整合時，不可直接於 dB 域做算術加總，而應先將各 beam contribution 轉為線性值後相加，最後再視需要轉換回 dB 表示。

\paragraph{(6) EPFD constraint}
為確保 NGSO--GSO coexistence 之法規合規性，系統必須滿足下列 aggregate EPFD 限制：
\begin{equation}
\mathrm{EPFD}_{\mathrm{agg}}(t)\le \mathrm{EPFD}_{\mathrm{lim}}.
\end{equation}

其中 $\mathrm{EPFD}_{\mathrm{lim}}$ 為對應場景與 reference bandwidth 下之 EPFD regulatory limit。

當某一 time slot 中 $\mathrm{EPFD}_{\mathrm{agg}}(t)>\mathrm{EPFD}_{\mathrm{lim}}$ 時，表示系統當前 beam power allocation 對受保護 GS 所造成之總干擾超出容許範圍，因此後續必須透過 power control 或其他資源調整機制降低 beam-level contributions，使系統重新滿足法規要求。

\subsection{Problem Formulation}

根據前述系統模型，本文考慮之資源分配問題同時涉及 user association、beam-level power allocation、co-channel interference，以及對受保護 GSO ground station 之 EPFD constraint。研究目標是在滿足 EPFD 法規限制之前提下，盡可能提升整體 user satisfaction。

令系統整體效益函數定義為所有 users satisfaction 之總和，如下所示。

\begin{equation}
\max_{\{x_{u,s,b}(t)\},\{P_{s,b}(t)\}}
\sum_{u\in\mathcal{U}(t)} \mathrm{sat}_u(t).
\end{equation}

則本文所考慮之聯合最佳化問題可表示如下。

\begin{equation}
\begin{aligned}
\max_{\{x_{u,s,b}(t)\},\{P_{s,b}(t)\}} \quad
& \sum_{u\in\mathcal{U}(t)} \mathrm{sat}_u(t)
\end{aligned}
\end{equation}

subject to
\begin{equation}
x_{u,s,b}(t)=0,
\qquad \forall u\notin\mathcal{C}_{s,b}(t),
\end{equation}

\begin{equation}
\sum_{s\in\mathcal{S}(t)}\sum_{b\in\mathcal{B}_s} x_{u,s,b}(t)\le 1,
\qquad \forall u\in\mathcal{U}(t),
\end{equation}

\begin{equation}
\sum_{b\in\mathcal{B}_s} P_{s,b}(t)\le P_s^{\max},
\qquad \forall s\in\mathcal{S}(t),
\end{equation}

\begin{equation}
0\le P_{s,b}(t)\le P_{s,b}^{\max},
\qquad \forall s\in\mathcal{S}(t),\ \forall b\in\mathcal{B}_s,
\end{equation}

\begin{equation}
\mathrm{EPFD}_{\mathrm{agg}}(t)\le \mathrm{EPFD}_{\mathrm{lim}},
\end{equation}

其中，其中，各限制式之物理意義說明如下。

式(27)如果 user u 不在衛星 s 的 beam b 覆蓋範圍內，那這個 association 不允許成立。

式(28)限制每位 user 於每個 time slot 中僅能被一個 satellite-beam 服務。

式(29)限制每顆衛星之總發射功率不得超過其 power budget。

式(30)限制各 beam 之發射功率不得超過 beam-level maximum power。

式(31)確保系統 aggregate EPFD 不超過 regulatory limit。

此目標函數用以在 EPFD 保護限制與衛星資源限制下，最大化可視 LEO 衛星系統中所有 users 的整體 satisfaction。需注意的是，由於 LEO 衛星具有時間變動之幾何關係，因此此最佳化問題需於每一個 time slot 重複求解。

\subsubsection{Discussion on Problem Complexity}

由上述式子可知，本文之最佳化問題屬於一個高複雜度之 mixed-integer non-linear optimization problem。特別是，beam power allocation 會同時影響 user-side SINR 與 GS-side aggregate EPFD，而 user reassociation 又會改變 beam loading，進一步影響 TDMA rate 與 satisfaction。因此，直接求得其全域最優解在實際計算上相當困難。

基於上述原因，本文於下一章提出一套分階段之 heuristic framework，在滿足 EPFD constraint 的前提下，逐步調整 beam powers 與 user reassociation，以提升整體 user satisfaction。